\scanpage{21}
\uline{Proposition} \deleted{\unsure{Pour tout|Remarque}} Deux arêtes orthogonales n'ont pas de sommet en commun.

En effet, les quatre arêtes orthogonales \unsure{à|:} $\{s,t\}$ sont $\{\pi_{s,u},\pi_{s,v}\}$, \deleted{et} $\{\pi_{t,u},\pi_{t,v}\}$ \deleted{et} d'une part, $\{u,v'\}$ et $\{v,u'\}$ d'autre part, elles ne rencontrent pas $\{s,t\}$

\uline{Corollaire} Pour tt $s\in S$, si $P(s)$ est l'ens des \deleted{\unsure{pts} $t\in S$ adjacents à $s$,} \inserted{arêtes incidentes à $s$}, l'application $a\mapsto\mathrm{tr}(\bar a)$ induit une bijection
\[
P(s)\longrightarrow T
\]
En effet, c'est une application injective entre deux ens de cardinal 5.

\uline{Corollaire} Soit $G_s^+$ le groupe des rotations de $\Pi$ fixant $s$ (i.e. les autom. fixant $s$, et induisant sur $\unsure{\mathscr P|P|\Phi}_{\geq s}$ une rotation \ldots) Alors \deleted{\unsure{$G_s^+$|$G_s$}} $T$ est un torseur sous $G_s^+$ -- donc si on fixe un gén. $u_s$ de $G_s^+$, $u_s$ est une permutation circulaire sur $T$

Comme $T$ est \inserted{de card} impair, on en conclut que $G_s^+$ opère par \uline{permutations paires}.

\uline{NB} Plus gén, on voit que les autom. d'un polygône combinatoire de rang \deleted{\unsure{impair sont des|impaire est dans}} $\equiv 1\ (4)$ (p.ex. d'un pentagône) induisent
% archive sheet number: 20
